\documentclass[11pt,journal,onecolumn]{IEEEtran}
\usepackage{amsmath,amssymb,amsfonts}
\usepackage{algorithmic}
\usepackage{graphicx}
\usepackage{textcomp}
\usepackage{xcolor}
\usepackage{hyperref}
\usepackage{booktabs}
\usepackage{bm}

\begin{document}

\title{An Inquiry into Robust Audio Fingerprinting: The Mathematical Foundations of the Shazam Search Algorithm}

\author{Aksel Aghajanyan \\ \small Technical Report --- Knowledge Base Archive}

\maketitle

\begin{abstract}
This paper provides a mathematically rigorous, self-contained exposition of the landmark audio search and fingerprinting algorithm proposed by Avery Li-Chun Wang in 2003, which serves as the core technology behind Shazam. We detail the entire pipeline: from discrete-time audio representation, the Short-Time Fourier Transform (STFT), and spectral peak extraction, to combinatorial hashing and time-coherent histogram-based matching. We analyze the mathematical properties that make this algorithm exceptionally robust to additive noise, MP3 compression artifacts, and frequency distortions, while scaling efficiently to databases containing millions of tracks.
\end{abstract}

\section{Introduction}
Identifying an unknown audio clip playing in a noisy environment (e.g., a cafe or a bar) is a challenging pattern recognition problem. Traditional hashing techniques (like MD5 or SHA-256) are fragile; changing a single bit in an audio stream completely alters the resulting hash. Therefore, an ``acoustic fingerprint'' must capture the perceptually salient, invariant features of the audio signal while discarding transient noise and channel distortions.

In 2003, Avery Li-Chun Wang published a breakthrough paper titled \textit{``An Industrial-Strength Audio Search Algorithm''}, introducing the algorithm underpinning the Shazam service. The key insight was to construct a sparse, noise-resistant representation of the audio signal based on local spectrogram maxima, paired combinatorially to form highly descriptive hashes. These hashes are mapped to a database, and matches are identified via a fast, linear time-offset alignment check.

This report unfolds the mathematical framework of this algorithm, providing proofs and formulas for each key stage.

---

\section{Mathematical Foundations of Audio Processing}

\subsection{Discrete-Time Audio Representation}
An analog audio signal $x_a(t)$ is digitized by sampling at a uniform rate $f_s = 1/T$, yielding the discrete-time sequence:
\begin{equation}
x[n] = x_a(nT), \quad n \in \mathbb{Z}
\end{equation}
Typical sampling rates are $f_s = 8000\text{ Hz}$ or $11025\text{ Hz}$ for telephony-quality analysis, up to $44100\text{ Hz}$ for high-fidelity audio.

\subsection{The Discrete Fourier Transform (DFT)}
For a finite discrete-time signal $x[n]$ of length $N$, its Discrete Fourier Transform $X[k]$ represents the signal in the frequency domain:
\begin{equation}
X[k] = \sum_{n=0}^{N-1} x[n] e^{-i \frac{2\pi}{N} k n}, \quad k \in \{0, 1, \dots, N-1\}
\end{equation}
where each bin $k$ corresponds to the physical frequency:
\begin{equation}
f_k = k \frac{f_s}{N}
\end{equation}

\subsection{Short-Time Fourier Transform (STFT) and Spectrogram}
Music is non-stationary; its frequency content changes over time. To capture time-varying spectral properties, we partition the signal $x[n]$ into overlapping blocks using a window function $w[n]$ (e.g., Hamming or Hanning window) of length $N$.

The Short-Time Fourier Transform (STFT) at time frame $m$ and frequency bin $k$ is defined as:
\begin{equation}
X(m, k) = \sum_{n=0}^{N-1} x[n + m H] w[n] e^{-i \frac{2\pi}{N} k n}
\end{equation}
where $H \in \mathbb{N}$ is the hop size (step size in samples between consecutive frames). The spectrogram $S(m, k)$ is the squared magnitude of the STFT, representing power spectral density:
\begin{equation}
S(m, k) = |X(m, k)|^2
\end{equation}
The spectrogram maps the continuous 1D audio waveform into a discrete 2D time-frequency grid $\mathbb{Z}^2$.

---

\section{Acoustic Fingerprinting: Peak Detection \& Constellation Map}

\subsection{Perceptual Saliency and Additive Noise Robustness}
Audio signals captured via mobile microphones are corrupted by ambient noise, reverberation, and frequency-dependent attenuation:
\begin{equation}
y[n] = (x * h)[n] + \eta[n]
\end{equation}
where $h[n]$ is the channel impulse response (reverberation, microphone EQ) and $\eta[n]$ is additive noise. In the frequency domain:
\begin{equation}
Y(m, k) \approx X(m, k) H(k) + \mathcal{N}(m, k)
\end{equation}
To achieve robustness, the algorithm extracts local peaks (spectral maxima) from the spectrogram $S(m, k)$. High-energy peaks survive additive noise ($\mathcal{N}(m, k)$ is overwhelmed by the signal energy at peak locations) and channel attenuation (which alters the absolute height of the peaks but generally preserves their positions as local maxima).

\subsection{Mathematical Definition of a Local Maximum}
A point $(m^*, k^*)$ is classified as a local maximum (or ``peak'') if its energy exceeds all neighbors within a surrounding neighborhood region $\mathcal{R}$:
\begin{equation}
S(m^*, k^*) > S(m, k) \quad \forall (m, k) \in \mathcal{R} \setminus \{(m^*, k^*)\}
\end{equation}
Typically, the neighborhood $\mathcal{R}$ is a bounding box or ellipse centered at $(m^*, k^*)$ spanning a size of $\Delta t_{\text{hood}}$ in time frames and $\Delta f_{\text{hood}}$ in frequency bins. For example:
\begin{equation}
\mathcal{R}(m^*, k^*) = \{ (m, k) \mid |m - m^*| \le D_t \text{ and } |k - k^*| \le D_f \}
\end{equation}
By tuning the dimensions $D_t$ and $D_f$, the algorithm controls the density of the extracted peaks. A typical target is to extract approximately $20$ to $30$ peaks per second. The resulting sparse representation of coordinates is the \textbf{Constellation Map} $C$:
\begin{equation}
C = \{ (t_i, f_i) \mid S(t_i, f_i) \text{ is a local maximum} \}
\end{equation}

---

\section{Combinatorial Hashing}

Directly matching individual peaks is highly prone to false positives because a single coordinate $(t_i, f_i)$ carries very little entropy (there are few frequency bins). To increase the specificity of each feature, Wang introduced \textbf{Combinatorial Hashing}, where multiple peaks are linked together.

\subsection{Anchor Points and Target Zones}
For each point $(t_a, f_a)$ in the constellation map designated as an \textbf{anchor point}, the algorithm identifies a \textbf{target zone} $\mathcal{Z}_a$ defined as a region shifted forward in time:
\begin{equation}
\mathcal{Z}_a = \{ (t, f) \in C \mid t_a + \delta_t \le t \le t_a + \delta_t + W_t, \text{ and } f_a - W_f \le f \le f_a + W_f \}
\end{equation}
where $\delta_t$ is a minimum delay (to avoid linking near-identical frequencies), $W_t$ is the width of the target zone in time, and $W_f$ is the half-frequency range.

For each point $(t_t, f_t) \in \mathcal{Z}_a$ (called a \textbf{target point}), a combinatorial fingerprint is generated.

\subsection{Hash Construction}
A single fingerprint hash $H_{j}$ consists of the tuple:
\begin{equation}
H_{j} = \Big( f_a, \, f_t, \, \Delta t \Big)
\end{equation}
where $\Delta t = t_t - t_a$. The offset time $t_a$ is recorded alongside the hash as the reference anchor time.

\begin{itemize}
    \item \textbf{Hash Key (32-bit integer):} Constructed by packing the bit-representations of $f_a$, $f_t$, and $\Delta t$.
    \item \textbf{Value mapping:} Each hash key is associated in the database with the tuple $(\text{Track ID}, t_a)$.
\end{itemize}

\subsection{Entropy and Combinatorial Scale}
Let $F$ be the number of frequency bins (e.g., $512$ bins for $N=1024$). Let the maximum time difference $\Delta t$ be bounded by $64$ frames.
If $f_a$ is encoded using 9 bits, $f_t$ using 9 bits, and $\Delta t$ using 6 bits, the hash key space requires $9 + 9 + 6 = 24$ bits ($16.7$ million unique combinations). This high entropy drastically minimizes accidental hash collisions between distinct tracks.

If an anchor point has on average $F_{\text{density}} \approx 10$ target points in its target zone, the database generates approximately $10 \times N_{\text{peaks}}$ hashes. This redundancy ensures that even if a large fraction (e.g., 80\%) of the spectral peaks are lost due to noise, the probability of at least a few shared hashes surviving remains extremely high.

---

\section{Database Matching and Time Coherence}

The search process takes a short query audio clip, extracts its constellation map, builds hashes, and queries the database.

\subsection{The Math of Relative Offsets}
Let a hash match occur between a query hash at anchor time $t_{a, q}$ and a database track entry at anchor time $t_{a, d}$ for a given $\text{Track ID}$.
If the query is a segment of the database track, then there must exist a constant temporal offset $t_{\text{offset}}$ such that:
\begin{equation}
t_{a, d} = t_{a, q} + t_{\text{offset}}
\end{equation}
We rewrite this relation as:
\begin{equation}
t_{\text{offset}} = t_{a, d} - t_{a, q}
\end{equation}
For all valid matching hashes belonging to the correct track, the difference $t_{a, d} - t_{a, q}$ will yield the exact same value $t_{\text{offset}}$.
For incorrect matches (spurious collisions), the differences $t_{a, d} - t_{a, q}$ will behave as randomly distributed noise variables.

\subsection{The Histogram Method}
Let $\mathcal{M}$ be the set of matching hash pairs returned from the database query:
\begin{equation}
\mathcal{M} = \Big\{ (d_i, t_{a, d}^{(i)}, t_{a, q}^{(i)}) \mid \text{Hash}_q^{(i)} = \text{Hash}_d^{(i)} \text{ and } d_i \text{ is the Track ID} \Big\}
\end{equation}
For each candidate track $d$, we define the set of offsets:
\begin{equation}
\mathcal{O}_d = \Big\{ t_{a, d}^{(i)} - t_{a, q}^{(i)} \mid (d_i, t_{a, d}^{(i)}, t_{a, q}^{(i)}) \in \mathcal{M} \text{ and } d_i = d \Big\}
\end{equation}
We compute the histogram of these offsets with a small bin size (e.g., corresponding to the frame resolution):
\begin{equation}
H_d(\tau) = \sum_{\Delta \in \mathcal{O}_d} \mathbb{I}(\Delta = \tau)
\end{equation}
where $\mathbb{I}(\cdot)$ is the indicator function. The matching score $S_d$ for track $d$ is the maximum peak in its offset histogram:
\begin{equation}
S_d = \max_{\tau} H_d(\tau)
\end{equation}
If $S_d$ exceeds a threshold $\theta$ (which can be statistically calibrated to yield near-zero false alarms), track $d$ is declared as the matching audio source.

\subsection{Visualizing Time Coherence}
Geometrically, if we plot $t_{a, d}$ versus $t_{a, q}$ for all matched hashes of the correct track, the true matches lie along a diagonal line of slope $1$:
\begin{equation}
t_{a, d} = 1 \cdot t_{a, q} + t_{\text{offset}}
\end{equation}
Spurious matches will be scattered randomly across the plane. The histogram search is mathematically equivalent to finding the line with the highest density of points parallel to the diagonal $y=x$, which is highly robust even when the signal-to-noise ratio (SNR) is extremely low.

---

\section{Performance and Complexity Analysis}

\subsection{Retrieval Complexity}
Let $N_q$ be the number of hashes generated from the query clip. The lookup time in the inverted index database for a single hash key is $O(1)$ on average. Thus, retrieving all matching candidate hits takes $O(N_q)$ database lookups.
Evaluating the histograms is linear in the number of retrieved hits. Since the hash keys have high entropy, the number of spurious hits is small, keeping the matching phase exceptionally fast.

\subsection{Robustness Analysis}
Let $p$ be the probability that a single spectral peak is correctly extracted in both the query and the database (under high noise).
The probability of a specific anchor-target pair surviving is $p^2$.
If we have $F_{\text{density}}$ target points per anchor, the expected number of matching hashes for a query containing $A$ anchors is:
\begin{equation}
\mathbb{E}[N_{\text{matching}}] = A \cdot F_{\text{density}} \cdot p^2
\end{equation}
Even if $p = 0.2$ (80\% of peaks are obliterated by noise), and we have a 10-second query with $A = 300$ anchors and $F_{\text{density}} = 10$:
\begin{equation}
\mathbb{E}[N_{\text{matching}}] = 300 \cdot 10 \cdot 0.04 = 120 \text{ matches}
\end{equation}
Since spurious matches are uniformly distributed across the track duration $T_{\text{track}}$, the probability of $K$ random matches landing in the exact same offset bin $\tau$ follows a Poisson distribution with parameter $\lambda$:
\begin{equation}
\lambda = \frac{N_{\text{spurious}}}{B}
\end{equation}
where $B$ is the number of histogram bins. Because $B$ is very large (e.g., a 3-minute track has $180 \text{ s} / 0.032 \text{ s} \approx 5600$ bins), $\lambda \ll 1$, making the probability of a false peak in the histogram extremely close to zero.

---

\section{Conclusion}
The Shazam audio fingerprinting algorithm is a masterclass in combining classical digital signal processing (DSP) with combinatorial heuristics and database indexation. By transforming the audio search problem into a time-coherent point-matching problem on a 2D spectrogram, it sidesteps the fragility of traditional audio comparison systems and delivers robust, sub-second search results under highly adverse acoustic conditions.

\begin{thebibliography}{1}
\bibitem{wang2003}
A. L.-C. Wang, ``An Industrial-Strength Audio Search Algorithm,'' in \textit{Proceedings of the International Conference on Music Information Retrieval (ISMIR)}, Baltimore, MD, 2003, pp. 7--13.
\end{thebibliography}

\end{document}